\section{相关工作}\label{sec:related}

\subsection{持续状态架构：状态与记忆}\label{sec:related-state}

本节把状态与记忆分开讲，因为两者是\textbf{实现与体现}的关系，不是并列的两个东西：
\textbf{状态实现记忆}——状态是前向过程中每一步实际读写的张量对象（实现层）；
\textbf{记忆是状态的宏观体现}——站在跨上下文的尺度看，「模型记得之前读过什么」这个能力，
就是状态被反复读写的累积效果。所以 \S\ref{sec:related-state} 先讲状态（微观、实现），再讲记忆（宏观、能力）。

\subsubsection{状态：架构的实现层}\label{sec:related-state-impl}

\textbf{定义。}状态是模型内部一组形状固定的张量对象，随前向逐步读写，不随历史长度增长。

\textbf{与 RNN 的关系：思想类似，实现差异很大。}两者都用固定大小的循环状态承载历史，
但本文研究的架构在四点上都不同于 RNN：
（i）不是单一隐向量的逐步递归，而是\textbf{四件套按时间尺度分工}
（delta 规则矩阵 memory + 逐层 c\_state + 滚动 queue + 中期 MTM，见 \S\ref{sec:model-state}）；
（ii）memory 的写入用 \textbf{delta 规则快权重}
\citep{schmidhuber1992learning,schlag2021linear,sun2023retnet,yang2024deltanet}，
而不是 RNN 的门控线性更新；
（iii）\textbf{规模不同}：RNN 的状态一般只有 KB--MB 级，本文的状态高出若干数量级（\S\ref{sec:bg}）——
「把容量做大」在参数规模上已有系统性的经验规律\citep{kaplan2020scaling}，
把同样的思路用在\textbf{状态}上正是本文的出发点；
（iv）\textbf{训练方式不同}：RNN/RWKV 训练时状态从零初始化、一批读完即清零，
单个状态最多跨一个样本（例如 2000 token）；本文的模型不按批次清空，
上一批结束时的状态直接续进下一批，整条训练流是一条全局连续的流（\S\ref{sec:train}）。
最后、也是最关键的一点：状态被设计成\textbf{可以整体读出、冻结、置换的显式对象}——
RNN 的隐状态没有对评测暴露这样的接口。

\textbf{与测试时训练/记忆模块一类工作的关系。}把「记忆」做成层内的可学习结构并非新想法：
快权重\citep{schmidhuber1992learning,ba2016using}、TTT\citep{sun2024ttt}、
Titans\citep{behrouz2025titans} 都让模型在前向中写入某种内部状态；
线性注意力/状态空间模型\citep{katharopoulos2020transformers,gu2023mamba,peng2023rwkv}
则把递归写成可并行扫描的形式。本文研究的架构在「读取用 delta 规则、写入用误差驱动外积」这一点上与它们同源，
区别在把该规则\textbf{嵌入到每一处线性投影}而不是只放在注意力位置，并把状态拆成四个时间尺度、
以显式对象的形式暴露给评测。本文的实验关心的不是「谁更强」，而是\textbf{这个状态里到底有什么、能带多远}。

\subsubsection{记忆：状态的宏观体现}\label{sec:related-mem}

\textbf{定义。}记忆 $=$ 状态在跨上下文信息传递上表现出的宏观能力；它与 \S\ref{sec:related-state-impl} 的状态
是同一个东西的两个观察尺度，\textbf{不是并列的两个模块}。

\textbf{与记忆增强路线的对照。}既有工作大多把记忆放在架构之外：外部记忆库或检索式增强
\citep{grave2017improving,lewis2020rag}从外部语料/存储里捞回信息；
滑窗或长窗口\citep{beltagy2020longformer,ding2023longnet}把记忆放进更长的窗口；
循环记忆 token\citep{bulatov2022recurrent}把记忆放在一串可训练槽里并通过反向传播更新。
这些路线的共同点是记忆\textbf{不在模型状态里}——要么在外存，要么在窗口，要么在参数。
本文研究的记忆完全由持续状态承载：不依赖重放、不借助外部存储、不靠额外参数。

\textbf{从「记忆全部落在持续状态上」到持续学习。}记忆既然全部落在持续状态上，
而状态在每次前向都被写入，那么模型在阅读过程中仍在改变自己的行为——这就是持续学习：
学习不只发生在训练时，也发生在推理的前向之中（与上下文内学习这一现象相关但并不相同：
ICL 是「看示例就会」\citep{brown2020gpt3,olsson2022induction,chan2022data}，
本文关心的是「读着读着状态变了、行为随之变了」）。这与「损失可塑性」\citep{dohare2024loss}是同一个问题的两个方向：既有工作关心\textbf{反复训练时会不会变坏}，
本文关心\textbf{一次前向流之内，行为会不会变好、变多少}。
可验证的含义是：在分布内/分布外的文本流上，模型的有效行为会随前向推进而改变——
本文用 A2/A6（写入的因果贡献与组件归因）、G 线（合成任务上的学习曲线）
以及 D 线（热启动记忆的领域特异性）三组证据来量这个「读着读着变了多少」。

\subsection{长上下文评测方法}\label{sec:related-eval}

\textbf{needle（needle-in-a-haystack）}\citep{kamradt2023needle}：
在长文本中埋入一条「针」，读完后提问，测模型能否精确检索。它是二值的（对/错），测的是\emph{检索}能力。
本文只把它作为边界预实验（\S\ref{sec:b4}）——固定大小状态的强项未必是精确定位，
所以如实限定 claim 边界，失败也写进 Limitations。

\textbf{滑窗基线}：用一个固定窗口宽度 $W$ 的滑动窗口模型作参照，
看固定状态的逐位置损失落在哪个 $W$ 上，从而把「状态记了多少」折算成一个直观数字「等效窗口宽度」（\S\ref{sec:b3}）。

\textbf{per-position 分析}：不只报一个总分，而是把损失\textbf{按位置逐点对齐}报告随位置的走向。
总分会掩盖结构：衰减形状、边界冲击、恢复长度只有逐位置曲线才看得见。
本文的位置衰减曲线（\S\ref{sec:b1}）与边界恢复曲线（\S\ref{sec:a3}）都属于这一口径。

\textbf{既有方法与本文的接点。}既有方法测的是「上下文（窗口）能看多远」，
本文测的是「固定状态能带多远」；两者底层指标口径相同（teacher-forced 逐位置损失），
因此结果可以互相对照——这也意味着本文的评测协议必须把「起点状态」写清楚（\S\ref{sec:cond}），
否则对照不成立。

\subsection{创新点总结}\label{sec:related-innovation}

\textbf{vs RNN}（\S\ref{sec:related-state-impl}）：思想类似——都用固定大小的循环状态承载历史；
但实现差异很大——不是单一隐向量的逐步递归，而是四件套按时间尺度分工，memory 用 delta 规则快权重写入；
规模上 RNN 状态只有 KB--MB 级，本文的状态高出若干数量级；训练方式上 RNN/RWKV 每批清零、
本文状态跨批次连续（\S\ref{sec:train}）；最关键的差别是状态被设计成\textbf{可整体读出、冻结、置换的显式对象}。

\textbf{vs 记忆增强路线 / 长上下文}（\S\ref{sec:related-mem}）：
既有工作把记忆放在架构之外——外部记忆库、检索增强、或单纯堆长上下文窗口；
本文的记忆完全由持续状态承载，不重放、不外挂，并在持续写入中得到持续学习。
一句话：别人的记忆在状态之外，本文的记忆就是状态在跨上下文尺度上的宏观体现。

\textbf{vs 既有评测方法}（\S\ref{sec:related-eval}）：
needle、滑窗基线、per-position 各自测的是「窗口能看多远」；本文沿用同一底层口径
（teacher-forced 逐位置损失，结果可相互对照），但把操纵旋钮从「改窗口」换成「动状态」——
热启动 / $\mathit{lock}$ / 重置 / 冻结全部落在状态接口上。

\textbf{vs 传统无状态架构}：状态即 KV cache，随输入即生即灭、不可操纵，
推理期不会因为读了新内容而改变自己；本文的状态是显式对象，可整体读出、冻结、置换，
且持续写入带来持续学习。这正对应贡献三的第一层表述——\textbf{状态是微观实现，可以显式调整}。

\textbf{汇总（围绕核心两点）}：本文的创新点 = ①\textbf{状态放大}（scaling up，
状态参数量相对 RNN 高出若干数量级，容量差两个到四个数量级）；
②\textbf{持续学习}（状态跨上下文持续写入，模型在前向中仍在改变自己的行为）。
这两点的可评测性由架构保证——状态是显式对象（可动），记忆可由位置衰减 / 重复距离 / 等效窗口三把尺量（可测）。
